% =====================================================================
%  3. Deriving gradient descent: Taylor, majorisation, gradient flow
% =====================================================================
\section{Three ways to arrive at gradient descent}

\begin{frame}{View 1: trust the linear model, but only nearby}
  The Taylor model at $\x_k$ is good close to $\x_k$ and gets worse as you move away. So minimise the model
  plus a penalty for going too far:
  \[
    \x_{k+1} = \argmin_{\x}\;\Bigl\{\underbrace{f(\x_k) + \nabla f(\x_k)\T(\x - \x_k)}_{\text{linear model}}
      + \underbrace{\tfrac{1}{2\alpha}\norm{\x - \x_k}^2}_{\text{stay close}}\Bigr\}.
  \]
  Set the gradient of the braces to zero:
  \[
    \nabla f(\x_k) + \tfrac1\alpha(\x - \x_k) = \mathbf 0
    \quad\Longrightarrow\quad
    \ans{\x_{k+1} = \x_k - \alpha\,\nabla f(\x_k)}
  \]
  \begin{itemize}\small
    \item $\alpha$ is the \alert{step size} (the ``learning rate'' in ML). A large $\alpha$ means we trust the linear model further out.
    \item Change the penalty and you get a different method. With $\tfrac12\norm{\cdot}_{\bP}^2$ you get preconditioned GD (Section~2).
  \end{itemize}
\end{frame}

\begin{frame}{View 2: minimise an upper bound}
  \begin{columns}[T]
    \begin{column}{0.47\textwidth}
      \small
      If $\nabla f$ is $L$-Lipschitz, the \emph{descent lemma} (Section~7, Problem M2) says that for every $\x$
      \[
        f(\x) \le f(\x_k) + \nabla f(\x_k)\T(\x-\x_k) + \tfrac L2\norm{\x-\x_k}^2 .
      \]
      \begin{itemize}\footnotesize
        \item That is View 1 with $\alpha = 1/L$, except now the model is an \alert{upper bound} touching $f$ at $\x_k$.
        \item Its minimiser is $\x_{k+1} = \x_k - \tfrac1L\nabla f(\x_k)$.
        \item $f(\x_{k+1}) \le$ bound at $\x_{k+1}$ $\le$ bound at $\x_k = f(\x_k)$, so $f$ \alert{has to go down}.
      \end{itemize}
      \footnotesize\textcolor{mGrey}{This trick is known as majorise and minimise (MM).}
    \end{column}
    \begin{column}{0.51\textwidth}
      \centering
      % 1-D descent lemma: f, tangent line, quadratic upper bound touching at x_k,
% the GD step as the minimiser of the upper bound, and the guaranteed drop.
% f(x) = 0.35 x^2 + 0.5 cos(2x) + 0.6  (f'' <= 0.7 + 2 = 2.7, so L = 2.7)
\begin{tikzpicture}[xscale=1.45, yscale=0.8, >={Stealth[length=5pt]},
  declare function={f(\x)=0.35*\x*\x + 0.5*cos(deg(2*\x)) + 0.6;
                    fp(\x)=0.7*\x - sin(deg(2*\x));
                    q(\x)=f(2.1) + fp(2.1)*(\x-2.1) + 0.5*2.7*(\x-2.1)*(\x-2.1);}]
  \def\xk{2.1}
  \def\Lc{2.7}
  \pgfmathsetmacro{\xn}{\xk - fp(\xk)/\Lc}
  \pgfmathsetmacro{\fk}{f(\xk)}
  \pgfmathsetmacro{\qn}{\fk - fp(\xk)*fp(\xk)/(2*\Lc)}
  \pgfmathsetmacro{\fn}{f(\xn)}
  % gap between the upper model and f
  \fill[mOrange!12, domain=0.7:2.75, samples=60]
    plot (\x, {q(\x)}) -- plot[domain=2.75:0.7, samples=60] (\x, {f(\x)}) -- cycle;
  % axes
  \draw[->, mGrey] (-0.6,0) -- (3.1,0) node[right, font=\scriptsize] {$x$};
  % models and f
  \draw[line width=1pt, mBlue, dashed, domain=1.4:2.8] plot (\x, {\fk + fp(\xk)*(\x-\xk)});
  \draw[line width=1.2pt, mOrange, domain=0.7:2.75, samples=60] plot (\x, {q(\x)});
  \draw[line width=1.8pt, mInk, domain=-0.4:2.9, samples=100] plot (\x, {f(\x)});
  % drop lines and points
  \draw[mGrey, densely dotted] (\xk,0) node[below, font=\scriptsize, mInk] {$x_k$} -- (\xk,\fk);
  \draw[mGrey, densely dotted] (\xn,0) node[below, font=\scriptsize, mInk] {$x_{k+1}$} -- (\xn,\qn);
  \draw[->, line width=1pt, mOrange] (\xk,-0.42) -- node[below=-1pt, font=\tiny] {$-\tfrac1L f'(x_k)$} (\xn,-0.42);
  % guaranteed decrease bracket
  \draw[decorate, decoration={brace, amplitude=3pt, mirror}, mGreen, line width=0.8pt]
    (\xn+0.06,\qn) -- (\xn+0.06,\fk) node[midway, right=2pt, font=\tiny, align=left] {drop $\ge$\\$\frac{1}{2L}f'(x_k)^2$};
  \draw[mGrey!70, densely dashed] (\xn,\fk) -- (\xk,\fk);
  \fill[mInk] (\xk,\fk) circle (1.7pt);
  \fill[mOrange] (\xn,\qn) circle (1.7pt);
  \fill[mGreen] (\xn,\fn) circle (1.7pt);
\end{tikzpicture}


      \scriptsize
      \textcolor{mInk}{\rule[2pt]{10pt}{1.8pt}} $f$ \quad
      \textcolor{mBlue}{\rule[2pt]{10pt}{1pt}} tangent \quad
      \textcolor{mOrange}{\rule[2pt]{10pt}{1.2pt}} upper model (curvature $L$)\\
      \textcolor{mOrange}{$\bullet$} minimum of the model, \quad \textcolor{mGreen}{$\bullet$} the true $f(x_{k+1})$, lower still.
    \end{column}
  \end{columns}
\end{frame}

\begin{frame}{View 3, the central idea: GD is forward Euler on the gradient flow}
  The \textbf{gradient flow} ODE $\;\dot\x(t) = -\nabla f(\x(t))$ always runs downhill:
  \[
    \frac{d}{dt}f(\x(t)) = \nabla f\T\dot\x = -\norm{\nabla f(\x(t))}^2 \le 0 .
  \]
  Now discretise it with \textbf{forward Euler}, step $h$: $\;\x_{k+1} = \x_k + h\,(-\nabla f(\x_k))$.
  That is \alert{gradient descent with $\alpha = h$}. Everything we know about Euler's method now applies.

  \vspace{4pt}
  \centering\small
  \begin{tabular}{@{}l l@{}}
    \toprule
    What we know from ODEs & What it means for optimisation \\
    \midrule
    Euler stability interval $|1 - h\lambda| < 1$ & the step limit $\alpha < 2/L$ (Section~5) \\
    stiffness (widely spread eigenvalues) & ill-conditioning and the zig-zag (Section~8) \\
    implicit (backward) Euler & the proximal-point method, stable for any step \\
    second-order ODE with friction & heavy-ball momentum (Section~11) \\
    $\ddot\x + \tfrac3t\dot\x + \nabla f = 0$ & Nesterov acceleration~\cite{su2016} \\
    \bottomrule
  \end{tabular}
\end{frame}

\statement{THE CENTRAL IDEA}{Gradient descent is forward Euler\\ on the gradient flow.}{So the step limit is a stability limit, and ill-conditioning is stiffness.}

\begin{frame}{The algorithm}
  \begin{columns}[T]
    \begin{column}{0.52\textwidth}
      \begin{thmbox}[Algorithm: gradient descent]
        \small
        \textbf{Input:} $\x_0$, step $\alpha>0$, tolerances $\varepsilon_{\text{rel}},\varepsilon_{\text{abs}}$, $k_{\max}$\\[2pt]
        \textbf{for} $k = 0,1,2,\dots,k_{\max}-1$:\\
        \quad $\g_k \gets \nabla f(\x_k)$\\
        \quad \textbf{if} $\norm{\g_k}\le\varepsilon_{\text{rel}}\norm{\g_0}+\varepsilon_{\text{abs}}$: \textbf{stop}\\
        \quad $\x_{k+1}\gets\x_k - \alpha\,\g_k$\\
        \textbf{return} $\x_k$
      \end{thmbox}
    \end{column}
    \begin{column}{0.46\textwidth}
      \small\textbf{When should it stop?}
      \begin{itemize}\footnotesize
        \item A small $\norm{\nabla f}$ is the natural test. For strongly convex $f$ it even bounds the error:
              $\norm{\x-\x^\star}\le\norm{\nabla f(\x)}/\mu$ (Section~6).
        \item Stopping only because $|f_{k+1}-f_k|$ is tiny is a trap. On a plateau $f$ hardly changes
              while $\x$ is still far from the answer (Problem A4).
        \item Always cap the iteration count and use a relative tolerance, because floating point limits
              how small the gradient can get (Section~9).
      \end{itemize}
    \end{column}
  \end{columns}
\end{frame}

\begin{frame}{One parameter by hand: the line fit, with numbers}
  \small
  Fix the slope at $m = 0.64$ and let GD find the intercept $b$, starting from $b_0 = 0$ with learning rate $\alpha = 0.1$.
  \[
    \frac{d\,\SSR}{db} = -2\sum_{i}\bigl(y_i - b - 0.64\,x_i\bigr),\quad
    \text{step} = \frac{d\,\SSR}{db}\times\alpha,\quad b_{\text{new}} = b_{\text{old}} - \text{step}.
  \]
  \centering\footnotesize
  \begin{tabular}{@{}c c c c c@{}}
    \toprule
    step $k$ & $b_k$ & slope $d\SSR/db$ & step size & $\SSR(b_{k+1})$ \\
    \midrule
    0 & $0$ & $-5.704$ & $-0.5704$ & $0.878$ \\
    1 & $0.5704$ & $-2.282$ & $-0.2282$ & $0.514$ \\
    2 & $0.7986$ & $-0.913$ & $-0.0913$ & $0.456$ \\
    3 & $0.8898$ & $-0.365$ & $-0.0365$ & $0.446$ \\
    \bottomrule
  \end{tabular}

  \raggedright
  \vspace{4pt}
  Each step is about $0.4\times$ the last (start $\SSR = 3.156$), and $b$ settles at \ans{b^\star = 0.951}.\par
  Steps shrink because the slope flattens near the minimum.
\end{frame}

\begin{frame}{Several parameters: the same recipe, once per parameter}
  \begin{columns}[T]
    \begin{column}{0.46\textwidth}
      \centering
      \scalebox{0.9}{% Gradient descent loop, steps (i)-(vi)
\begin{tikzpicture}[
  node distance=0.30cm,
  step/.style={draw=#1, line width=0.8pt, rounded corners=4pt, fill=white,
               drop shadow={opacity=0.15, shadow xshift=0.5pt, shadow yshift=-0.7pt},
               text width=3.8cm, minimum height=0.62cm, align=left, font=\scriptsize\color{mInk},
               inner xsep=10pt},
  step/.default=mNavy,
  badge/.style={circle, fill=#1, text=white, font=\tiny\bfseries, inner sep=1pt, minimum size=10pt},
  arr/.style={-{Stealth[length=5pt]}, line width=0.9pt, mGrey}]
  \node[step] (s1) {take the derivative of the loss for each parameter};
  \node[step, below=of s1] (s2) {pick a starting value for each parameter};
  \node[step, below=of s2] (s3) {plug the current values into the derivatives};
  \node[step=mOrange, below=of s3] (s4) {step size $=$ slope $\times$ learning rate};
  \node[step=mOrange, below=of s4] (s5) {new value $=$ old value $-$ step size};
  \node[step=mTeal, below=of s5] (s6) {stop if the step is tiny or max steps reached};
  \foreach \a/\b in {s1/s2, s2/s3, s3/s4, s4/s5, s5/s6} \draw[arr] (\a) -- (\b);
  \draw[arr] (s6.east) -- ++(0.55,0) |- node[pos=0.25, right=3pt, font=\scriptsize, mGrey] {not yet} (s3.east);
  \foreach \n/\lab/\c in {s1/i/mNavy, s2/ii/mNavy, s3/iii/mNavy, s4/iv/mOrange, s5/v/mOrange, s6/vi/mTeal}
    \node[badge=\c] at (\n.west) {\lab};
\end{tikzpicture}
}
    \end{column}
    \begin{column}{0.52\textwidth}
      \small
      For $\hat y = b + mx$ both parameters move:
      \begin{align*}
        \frac{\partial\SSR}{\partial b} &= -2\sum_i\bigl(y_i - b - m x_i\bigr),\\
        \frac{\partial\SSR}{\partial m} &= -2\sum_i x_i\bigl(y_i - b - m x_i\bigr).
      \end{align*}
      In general, for every $w_j$ at once:
      \[
        \ans{w_j \leftarrow w_j - \alpha\,\frac{\partial f}{\partial w_j}}
      \]
      \textbf{Stop} when the step size is near zero (say $<10^{-3}$) or after a maximum number of steps.
    \end{column}
  \end{columns}
\end{frame}

\begin{frame}[fragile]{The whole method in a dozen lines}
  \begin{columns}[T]
    \begin{column}{0.55\textwidth}
\begin{lstlisting}[style=py]
import numpy as np

def gradient_descent(grad, x0, alpha,
                     rtol=1e-8, maxit=10_000):
    x = np.asarray(x0, dtype=float).copy()
    g0 = np.linalg.norm(grad(x))
    for k in range(maxit):
        g = grad(x)
        if np.linalg.norm(g) <= rtol * g0:
            break
        x -= alpha * g       # O(n) update
    return x, k

# least squares: f(w) = ||Xw - y||^2 / (2m)
def ls_grad(w, X, y):
    return X.T @ (X @ w - y) / len(y)
\end{lstlisting}
    \end{column}
    \begin{column}{0.43\textwidth}
      \small\textbf{What one iteration costs}
      \begin{itemize}\footnotesize
        \item One gradient and an $O(n)$ vector update, with $O(n)$ memory.
        \item For least squares with $\mathbf X\in\R^{m\times n}$, that is two matrix--vector products, $O(mn)$.
              Nothing gets factorised.
        \item Written with array operations (BLAS, GPUs) instead of Python loops, the same code runs for
              $n = 2$ or for millions of variables.
        \item Total cost is (number of iterations) $\times$ (cost of $\nabla f$). The rest of the module is
              about that first factor.
      \end{itemize}
    \end{column}
  \end{columns}
\end{frame}
